\documentclass[11pt]{article}

\usepackage[margin=1in]{geometry}
\usepackage{amsmath,amssymb,amsthm}
\usepackage{booktabs}
\usepackage{hyperref}
\usepackage{url}
\usepackage{graphicx}
\usepackage{siunitx}

\hypersetup{colorlinks=true, linkcolor=blue, citecolor=blue, urlcolor=blue}

\newtheorem{proposition}{Proposition}
\newtheorem{corollary}[proposition]{Corollary}
\theoremstyle{definition}
\newtheorem{definition}[proposition]{Definition}
\theoremstyle{remark}
\newtheorem{remark}[proposition]{Remark}

\newcommand{\dv}{\Delta v}
\newcommand{\R}{\mathbb{R}}
\newcommand{\norm}[1]{\lVert #1 \rVert}

\title{Certified Fleet-Wide Collision Avoidance:\\
Constraining Induced Conjunctions Inside the Optimizer,\\
and Measuring What It Costs}

\author{AEGIS Project\thanks{Source, contracts, and every number in this paper:
\url{https://github.com/} --- see \texttt{aegis/specs/step14\_fleet\_optimization.md}
for the implementation contract and \texttt{docs/} for the measurement record.}}

\date{September 2026}

\begin{document}
\maketitle

\begin{abstract}
Collision avoidance maneuver (CAM) optimization is mature for one spacecraft
facing passive secondaries. It is not mature for a fleet, and the gap is
specific: the multi-encounter literature assumes that resolving one conjunction
does not create another, while measurements at mega-constellation scale show
that \SI{81.4}{\percent} of Starlink-on-Starlink maneuvers are cascade-induced.
Operators handle the cascade with an outer re-screening loop, which can
discover an induced conjunction but cannot prevent one, and offers no
completeness guarantee.

We put the condition inside the optimizer. Three results make that tractable
and trustworthy. First, a \emph{reachability bound}: given a per-satellite
$\dv$ budget, an inflated, time-varying screening gate provably enumerates
every pair that \emph{any} admissible plan could bring within an exclusion
radius, converting ``re-screen and hope'' into a certified candidate set.
Second, an \emph{exact-penalty threshold}: the apparent $\dv$-versus-risk
tradeoff observed under a soft penalty is an artifact of choosing the penalty
below the largest optimal dual, and above that computable threshold the soft
and hard formulations coincide. Third, the $\dv$-versus-induced-risk Pareto
frontier is \emph{convex and piecewise linear} with dual-valued slopes, so it
is computable exactly rather than sampled; we verify convexity to
\num{9e-19} and slope--dual agreement to \num{1.5e-13}.

We measure the \emph{safety premium} --- the extra $\dv$ a coordinated plan
pays to create no new conjunction --- and find it is governed by how tightly
the fleet is packed rather than by the number of conjunctions: a median of
\SI{2.7}{\percent} with a p90 of \SI{8.1}{\percent} for in-plane neighbour
spacings of \SIrange{4.5}{12}{\kilo\meter}, rising to infeasibility at
\SI{3}{\kilo\meter}. The median sits inside the \SIrange{1.1}{3.8}{\percent}
range published for single-spacecraft multi-encounter CAM under stronger
probability margins, from an entirely different formulation.

Finally we report a certified acceleration and a negative result. Solving with
a subset of the induced-conjunction rows and closing the gap with a
lazy-constraint loop returns the \emph{exact} optimum of the full problem for
any guess --- verified to machine epsilon over an empty, a learned, and an
adversarial guess --- and yields \numrange{1.8}{9.0}$\times$ speedups as the
problem grows. A physics-informed graph neural network trained to supply that
guess did \emph{not} beat starting from the empty set, for reasons we quantify.
\end{abstract}

\section{Introduction}

The object of study is one sentence: given a network of simultaneous
conjunctions across a maneuverable fleet, compute one coordinated set of
impulsive maneuvers minimizing total fleet $\dv$ subject to resolving every
modelled conjunction and creating no new one.

The single-spacecraft version of this problem is well solved. Convex and
sequential-convex formulations compute fuel-optimal CAMs against multiple
simultaneous encounters~\cite{armellin2021,pavanello2024}, with keep-out zones
convexified by the supporting-hyperplane technique introduced for ellipsoidal
zones by Mueller~\cite{mueller2009} and formalized as projection-and-lin\-ear\-i\-za\-tion
by Mao et al.~\cite{mao2017}. That literature is explicit about what it
assumes. Pavanello et al.\ state it plainly: ``from an operational point of
view, there is no reason to suppose that a close approach with a particular
secondary is likely to promote a close approach with some other
secondary''~\cite{pavanello2024}.

Chen et al.\ measured that supposition at scale and found it
false~\cite{chen2024}: \SI{81.4}{\percent} of Starlink-on-Starlink maneuvers
are cascade-induced, a single external avoidance event triggered as many as 41
induced maneuvers, and \SI{79.5}{\percent} of network lifetime is lost to
cascades. The gap between those two sentences is this paper.

Operationally the cascade is handled by an outer loop: generate the
post-maneuver ephemeris, resubmit it for a fresh screening, repeat. NASA's
Conjunction Assessment Best Practices Handbook describes exactly this and calls
it the best available approach~\cite{nasa2023}. The loop has two structural
limits. It has no termination guarantee, and --- more importantly --- it
re-screens the one trajectory the optimizer happened to choose, so it can only
discover conjunctions that plan created. It cannot say which pairs the
optimizer had the \emph{power} to endanger, which is exactly what a constraint
needs to know.

\subsection{Contributions}

\begin{enumerate}
\item \textbf{A reachability-certified candidate set} (Propositions~\ref{prop:reach}
and~\ref{prop:complete}). A bound on displacement under a $\dv$ budget yields
an inflated screening gate that provably enumerates every pair any admissible
plan could turn into a conjunction. We find no prior art for a completeness
proof over a \emph{set} of admissible-but-unknown maneuvers; the closest
work~\cite{rivero2025} requires the actual planned ephemeris and establishes
zero false negatives empirically.

\item \textbf{Induced-conjunction constraints inside the program}, with a
second-order inflation derived from the \emph{tidal} relative acceleration
rather than from relative speed. The distinction is not cosmetic: at a
\SI{30}{\second} step the na\"ive Lipschitz bound is \SI{285}{\kilo\meter} and
the tidal bound is \SI{0.55}{\kilo\meter}, and only the second leaves a
\SI{2}{\kilo\meter} exclusion radius meaning anything.

\item \textbf{Structural results on the tradeoff} (Propositions~\ref{prop:penalty}
and~\ref{prop:frontier}): an exact-penalty threshold that dissolves the apparent
fuel-versus-risk tradeoff, and convex piecewise-linear structure that makes the
Pareto frontier exactly computable. Both are classical optimization theorems;
the contribution is the instantiation and the verification.

\item \textbf{A Farkas certificate} naming which conjunctions are mutually
irreconcilable when no zero-induced plan exists, so ``infeasible'' becomes
``resolving $A$ and $B$ together necessarily brings $C$ inside the exclusion
radius.''

\item \textbf{The safety premium, measured} and stratified by graph structure
and by fleet packing --- to our knowledge the first fleet-level, $\dv$-denominated
cost-of-safety curve for induced conjunctions.

\item \textbf{A certified acceleration, and an honest negative result} on the
learned component (Section~\ref{sec:accel}).

\item \textbf{A relative-velocity condition on the inflation bound.} The
second-order inflation of contribution~2 bounds the \emph{curvature} of the
relative trajectory and is therefore conditional on the relative velocity at
the constrained minimum being small --- a condition that holds for co-orbital
pairs, fails for crossings, and is not usually stated. Where it fails a
sampled grid locates nothing: we exhibit a \SI{7.24}{\kilo\meter\per\second}
crossing whose \SI{143}{\meter} minimum lies between \SI{30}{\second} samples
reading \SI{76}{} and \SI{141}{\kilo\meter}. Bracketing by range-rate sign
change and refining by cubic Hermite restores the guarantee independently of
the step, and is what makes the zero result hold across budgets rather than at
one.

\item \textbf{Lazy epoch generation, and the two failure modes it exposes.}
Constraining local minima of the \emph{nominal} separation is insufficient ---
for a co-orbital pair that separation is flat, and the maneuver itself creates
the time-variation --- so the epochs that matter are discovered from the solved
plan rather than guessed beforehand. Once rows are generated this way, two
things break that do not break otherwise: a formulation that pins an attained
safety level measured on the enumerated rows becomes infeasible the moment a
row is added and silently falls back to an unrefined plan, and a certificate
evaluated on the enumerated rows will call a plan safe while it violates a row
the generation itself produced. Both are stated and fixed in
Section~\ref{sec:experiments}.
\end{enumerate}

We are explicit about what is \emph{not} new. The supporting-hyperplane
convexification is standard~\cite{mueller2009,mao2017,armellin2021} and is used
here as machinery, not claimed. The exact-penalty and parametric-LP theorems
are textbook~\cite{bertsekas1999,schrijver1986}.

\section{Formulation}

\subsection{Decision variables and displacement}

Let $\mathcal{F}$ be the maneuverable fleet. Satellite $i$ has candidate burn
epochs $\tau_{i,1} < \cdots < \tau_{i,K_i}$ and impulses
$\dv_{i,k} \in \R^3$ in its RTN frame, stacked into
$x \in \R^{3\sum_i K_i}$.

Everything the optimizer knows about dynamics enters through the exact
Clohessy--Wiltshire position-from-impulse block, with $s=\sin(n\sigma)$,
$c=\cos(n\sigma)$:
\begin{equation}
\Phi(n,\sigma) =
\begin{bmatrix}
s/n & 2(1-c)/n & 0\\
2(c-1)/n & (4s-3n\sigma)/n & 0\\
0 & 0 & s/n
\end{bmatrix}.
\label{eq:phi}
\end{equation}
The $(2,2)$ entry carries the secular term $-3\sigma$: only a tangential
impulse buys displacement growing without bound in lead time, while radial and
cross-track impulses produce bounded periodic motion of amplitude at most
$2\dv/n$. That asymmetry is why lead time, not propellant, is the scarce
resource.

The displacement of satellite $i$ at time $t$, in ECI, is linear in $x$:
\begin{equation}
\delta r_i(t) = \Psi_i(t)\,x,
\qquad
\Psi_i(t) = R_i(t)\big[\,\cdots\ \Phi(n_i, t-\tau_{i,k})\,\mathbf{1}[\tau_{i,k}\le t]\ \cdots\,\big],
\end{equation}
with $R_i(t)$ the RTN-to-ECI rotation. Rotating through each satellite's own
frame --- rather than approximating the frame difference with a scalar
alignment term --- is what makes the following correct for crossing geometries;
Section~\ref{sec:fidelity} quantifies what the approximation costs.

\subsection{B-plane miss sensitivity}

For conjunction $j$ between objects $a,b$ at nominal TCA $t_j$, write
$d_j = r_b - r_a$, $w_j = v_b - v_a$, and let
$P_j = I_3 - \hat w_j \hat w_j^{\!\top}$ be the projector onto the encounter
B-plane. Let $B_j = \Psi_b(t_j) - \Psi_a(t_j)$.

\begin{proposition}[First-order miss distance]
\label{prop:miss}
The post-maneuver minimum separation near $t_j$ is
\begin{equation}
m_j(x) = \norm{P_j\,(d_j + B_j x)} + O(\norm{\cdot}^2).
\end{equation}
\end{proposition}

\begin{proof}
Near TCA the relative motion is $d(t) = d_j + B_j x + w_j (t-t_j) + O((t-t_j)^2)$.
Minimizing the norm over $t$ leaves the residual $\norm{P_j(d_j + B_jx)}$; the
neglected relative acceleration acts over a time shift that is itself first
order in the displacement, so its contribution is second order.
\end{proof}

The projector matters because a displacement component along $w_j$ changes the
\emph{time} of closest approach, not its \emph{distance}. Omitting it credits a
burn with separation it does not buy.

The safety requirement $m_j(x) \ge \rho_j$ is reverse-convex. We use the
standard supporting-hyperplane restriction: for any unit $u$,
\begin{equation}
u^{\!\top} P_j (d_j + B_j x) \ \ge\ \rho_j
\quad\Longrightarrow\quad
\norm{P_j(d_j + B_j x)} \ \ge\ \rho_j
\label{eq:restriction}
\end{equation}
by Cauchy--Schwarz. The linear constraint is therefore an \emph{inner}
approximation for \emph{any} $u$ --- a fact we lean on heavily in
Section~\ref{sec:accel}, because it means no direction supplier, learned or
otherwise, can make a plan unsafe.

\begin{proposition}[Monotone refinement]
\label{prop:scp}
If $x^\star$ solves the program linearized at directions $u$, and
$u_j' = P_j(d_j + B_jx^\star)/\norm{P_j(d_j+B_jx^\star)}$, then $x^\star$
remains feasible for the program linearized at $u'$. Hence the optimal cost is
non-increasing across refinements and, bounded below, converges.
\end{proposition}

\begin{proof}
$u_j'^{\!\top}P_j(d_j+B_jx^\star) = \norm{P_j(d_j+B_jx^\star)} \ge \rho_j$ by
\eqref{eq:restriction} applied at $u$.
\end{proof}

\section{Induced conjunctions}

\subsection{Latent constraints}

Let $\mathcal{L}$ index \emph{latent} pair/epoch combinations --- pairs with at
least one maneuverable member that are not currently conjunctions. For each we
impose $\norm{d_p + B_p x} \ge s_{\min} + \gamma$, restricted as in
\eqref{eq:restriction}.

The inflation $\gamma$ covers what a first-order model evaluated at a grid
sample does not see:
\begin{equation}
\gamma = \tfrac12 A h^2 + \Lambda h,
\qquad
A \le \frac{2\mu G}{r_{\min}^{3}},
\qquad
\Lambda = \sqrt{51}\,(D_a + D_b),
\label{eq:tidal}
\end{equation}
where $h$ is the grid step, $G$ the gate, and $D_i$ the per-satellite budget.
$A$ bounds the \emph{tidal} relative acceleration of a pair already inside the
gate --- two objects a hundred kilometres apart feel nearly the same gravity,
and only the residual bends their relative trajectory. $\Lambda$ bounds how
fast the maneuver-induced displacement itself changes, via the Lipschitz
constant of $\Phi$.

\begin{remark}
The alternative --- bounding the change between samples by the relative speed,
$\tfrac12 v_{\max} h$ with $v_{\max} = \SI{19}{\kilo\meter\per\second}$ ---
gives \SI{285}{\kilo\meter} at $h = \SI{30}{\second}$, which swamps any
realistic exclusion radius and makes the constraint set vacuous. Equation
\eqref{eq:tidal} gives \SI{0.55}{\kilo\meter} at a \SI{1}{\meter\per\second}
budget and \SI{0.14}{\kilo\meter} at \SI{50}{\milli\meter\per\second}. The
factor of \num{500} to \num{2000} is the difference between a usable
formulation and an empty one.
\end{remark}

\subsection{Reachability and completeness}

\begin{proposition}[Displacement bound]
\label{prop:reach}
Let satellite $i$ carry total impulse budget $D_i = \sum_k \norm{\dv_{i,k}}$
with earliest burn epoch $\tau_i^{\min}$. Then for every admissible plan and
every $t$,
\begin{equation}
\norm{\delta r_i(t)} \;\le\; D_i \max_{\tau \in [\tau_i^{\min},\,t]} \norm{\Phi(n_i, t-\tau)}_2
\;=:\; \rho_i(t).
\end{equation}
\end{proposition}

\begin{proof}
Triangle inequality and submultiplicativity of the spectral norm.
\end{proof}

Evaluating the maximum requires care: $\norm{\Phi(n,\sigma)}_2$ is \emph{not}
monotone in $\sigma$. Its growth rate is $3 - 4\cos(n\sigma)$, negative
whenever $\cos(n\sigma) > 3/4$, so the norm dips just after every whole orbit
of lead. We therefore use the closed-form monotone envelope
\begin{equation}
\norm{\Phi(n,\sigma)}_2 \le \norm{\Phi(n,\sigma)}_F \le
\frac{\sqrt{(3n\sigma+4)^2 + 34}}{n} =: M(n,\sigma),
\label{eq:envelope}
\end{equation}
which follows from $2\sin^2 \le 2$, $8(1-\cos)^2 \le 32$ and
$|4\sin(n\sigma) - 3n\sigma| \le 4 + 3n\sigma$. $M$ is asymptotic to $3\sigma$
and over-estimates by at most about $7/n$ seconds.

\begin{proposition}[Candidate-set completeness]
\label{prop:complete}
If a pair $(a,b)$ satisfies
$\norm{d_{ab}(t)} > s_{\min} + \rho_a(t) + \rho_b(t)$ for every $t$ in the
window, then no admissible plan can bring it within $s_{\min}$. Screening the
\emph{nominal} catalog with that inflated, time-varying gate therefore
enumerates every pair any admissible plan could turn into a conjunction.
\end{proposition}

\begin{proof}
$\norm{d_{ab}(t) + \delta r_b(t) - \delta r_a(t)} \ge \norm{d_{ab}(t)} -
\rho_a(t) - \rho_b(t) > s_{\min}$.
\end{proof}

\begin{remark}
The gate is linear in the budget, and asymptotically $\rho_i(t) \approx
3 D_i (t - \tau_i^{\min})$ --- the secular along-track term and nothing else.
At \SI{1}{\meter\per\second} over one day that is \SI{263}{\kilo\meter} per
satellite; at the \SI{5}{\centi\meter\per\second} avoidance actually uses, it
is \SI{13}{\kilo\meter}, comparable to the \SI{17}{\kilo\meter} along-track
half-width of NASA CARA's operational LEO volume. Budgeting the gate is
therefore the same act as budgeting the fuel.
\end{remark}

\section{Structure of the solution}

With mandatory slack $\sigma$ on every safety row --- required, because a
purely radial or cross-track geometry has bounded achievable response and no
finite $\dv$ satisfies it --- the program is
\begin{equation}
\min_{x,\,s,\,\sigma}\ \ c^{\!\top}x + \lambda \mathbf{1}^{\!\top}s + \mu \mathbf{1}^{\!\top}\sigma
\quad \text{s.t.} \quad
\begin{cases}
u_j^{\!\top}P_j(d_j + B_jx) + s_j \ge \rho_j, & j \in \mathcal{J}\\
\hat u_p^{\!\top}(d_p + B_px) + \sigma_p \ge f_p, & p \in \mathcal{L}\\
\mathbf{1}^{\!\top}\sigma \le \varepsilon, & \\
x \in X, \ s,\sigma \ge 0. &
\end{cases}
\label{eq:program}
\end{equation}

\begin{proposition}[Exact penalty]
\label{prop:penalty}
For the linear program $\min\{c^{\!\top}x : Ax \ge b,\ x \in X\}$ with optimal
dual $y^\star$, every optimal solution of its penalized form has $\sigma = 0$
whenever the hard problem is feasible and $\lambda > \norm{y^\star}_\infty$.
\end{proposition}

\begin{proof}
The dual of the penalized problem is the dual of the hard problem with the
added box $y \le \lambda\mathbf{1}$. If $\lambda > \norm{y^\star}_\infty$ then
$y^\star$ is interior to that box, so the dual --- and hence the primal ---
optimum is unchanged; and any $\sigma>0$ costs $\lambda$ per unit while buying
at most $\norm{y^\star}_\infty < \lambda$ of relief.
\end{proof}

The consequence for this domain is the point: \emph{the apparent
$\dv$-versus-risk tradeoff observed under a soft penalty is an artifact of
choosing $\lambda \le \norm{y^\star}_\infty$}, not a property of the physics.
Above the computable threshold the formulations coincide.

\begin{proposition}[Frontier structure]
\label{prop:frontier}
$V^\star(\varepsilon)$, the optimal cost of \eqref{eq:program} as a function of
the induced-risk budget, is convex, piecewise linear and non-increasing in
$\varepsilon$, with finitely many breakpoints, and its left derivative equals
minus the optimal multiplier of the budget row.
\end{proposition}

\begin{proof}
Standard right-hand-side sensitivity: $V^\star$ is the support function of the
dual feasible set, a pointwise maximum of affine functions of $\varepsilon$,
over a polyhedron with finitely many vertices.
\end{proof}

The domain consequence is that the $\dv$-versus-induced-risk Pareto frontier is
\emph{exactly} representable by finitely many breakpoints and computable by
parametric LP, rather than sampled as a point cloud --- and that convexity is a
checkable invariant rather than an impression.

Finally, when no zero-induced plan exists, a deletion filter over the safety
rows returns a minimal infeasible subsystem, which renders in English as
``resolving $A$ and $B$ together necessarily brings $C$ within $s_{\min}$.''
Farkas' lemma is the basis~\cite{schrijver1986}; the contribution is the
mapping back to named conjunctions.

\section{Certified acceleration}
\label{sec:accel}

Solve \eqref{eq:program} with a \emph{subset} of $\mathcal{L}$, evaluate every
omitted row at the solution, add back any that is violated, and repeat.

\begin{proposition}[Lazy exactness]
\label{prop:lazy}
At termination the point is feasible for the full program, and its cost is a
lower bound on the full optimum because every intermediate program is a
relaxation. Hence the returned solution is optimal for the full problem
whatever subset was guessed.
\end{proposition}

\begin{remark}[A trap]
The omitted rows must be checked with the \emph{linearized} residual
$u^{\!\top}(d + Bx) - f$, not the physical separation $\norm{d+Bx} - f$. By
\eqref{eq:restriction} the linear form is never larger, so a point can satisfy
the true separation while violating the row the program enforces; checking the
norm terminates the loop early on a solution the full program would have
rejected. Using the physical residual broke exactness on 6 of 27 scenarios in
our implementation before the distinction was made explicit.
\end{remark}

Proposition~\ref{prop:lazy} is what makes a learned prior safe to use: the
prediction chooses only which rows to try first, and the loop certifies the
answer. We train a physics-informed graph network (\num{892}k parameters) over
the conjunction graph to predict the active set and dual prices, with a loss
that pushes the predicted allocation through the stored sensitivity tensors to
evaluate real constraint residuals. The result is reported in
Section~\ref{sec:results-accel}, including where it fails.

\section{Experiments}
\label{sec:experiments}

All scenarios are deterministic and seeded; every probability inherits
TLE-grade covariance and is labelled as such. Ground truth throughout is SGP4:
a plan is applied to the catalog, both catalogs are propagated, and the
realized geometry is measured.

\subsection{Model fidelity}
\label{sec:fidelity}

Impulses are applied to SGP4 objects by a first-order Gauss variational update
followed by a mean-element refit against the exact post-burn state. Measured
over 160 samples (40 mixed-axis impulses spanning
\SIrange{1}{500}{\milli\meter\per\second}, each at four lead times):

\begin{center}
\begin{tabular}{lrrr}
\toprule
Impulse magnitude & $n$ & refit converged & median relative error \\
\midrule
$\ge$ \SI{20}{\milli\meter\per\second} & 88 & \SI{100}{\percent} & \SIrange{0.22}{0.24}{\percent} \\
$<$ \SI{20}{\milli\meter\per\second} & 72 & \SI{17}{\percent} & \SI{11.6}{\percent} \\
all & 160 & \SI{62}{\percent} & \SI{0.28}{\percent} \\
\bottomrule
\end{tabular}
\end{center}

Per-axis, for converged fits: transverse \SI{0.20}{\percent} median
(\SI{0.42}{\percent} max), radial \SI{0.43}{\percent} (\SI{4.2}{\percent}),
normal \SI{2.6}{\percent} (\SI{7.6}{\percent}). Cross-track error grows with
lead time --- differential $J_2$ nodal regression from the inclination change,
which CW does not model --- which is why the optimizer prices cross-track burns
at five times along-track. The absolute disagreement is about ten metres
against an exclusion radius of \SI{2}{\kilo\meter}.

\subsection{The scalar sensitivity model fails on intra-fleet geometry}

A scalar along-track model estimates the secondary's contribution as
$(\hat d \cdot \hat t_a)(\hat t_a \cdot \hat t_b)$ where the correct
coefficient is $\hat d \cdot \hat t_b$. Those agree only when the orbit planes
nearly coincide. Measured on cross-plane intra-fleet conjunctions where the
\emph{secondary} maneuvers (24 cases, \SI{106}{\degree} plane separation,
\SI{12.1}{\kilo\meter\per\second} relative speed):

\begin{center}
\begin{tabular}{lrrr}
\toprule
$\dv$ & true change & B-plane model & scalar model \\
\midrule
\SI{10}{\milli\meter\per\second} & \SI{-0.1450}{\kilo\meter} & \SI{-0.1473}{\kilo\meter} & \SI{+0.0424}{\kilo\meter} \\
\SI{50}{\milli\meter\per\second} & \SI{-0.6660}{\kilo\meter} & \SI{-0.7317}{\kilo\meter} & \SI{+0.2119}{\kilo\meter} \\
\bottomrule
\end{tabular}
\end{center}

The scalar model has the \emph{wrong sign}: it predicts the burn increases
separation when it decreases it, with a median error of \SI{128}{\percent}
against \SI{1.0}{\percent} for the B-plane form. Co-planar pairs in the same
experiment agree to \SI{0.01}{\percent}, exactly as predicted. For coordinated
multi-satellite maneuvering this is a correctness fix, not a refinement.

\subsection{The safety premium}

Eighty deterministic \texttt{induced-cascade} scenarios: one maneuverable
satellite forced to act by an external hazard, with in-plane fleet-mates close
enough that the avoidance displacement can reach them. Neighbour spacing is
sampled per seed over \SIrange{3}{15}{\kilo\meter}, so the seed sweep is also a
sweep over how tightly the fleet is packed. The exclusion radius is derived
from the assessed covariances rather than chosen --- the separation at which
$P_c$ reaches the watch threshold, times a \num{1.5} margin, median
\SI{3.137}{\kilo\meter}.

Induced conjunctions are counted twice. \emph{Predicted} is the optimizer's own
constraint residual; \emph{measured} is a full SGP4 re-screen of the maneuvered
catalog, counting pairs that reach WATCH and were not there before. Only the
second is evidence.

\begin{center}
\begin{tabular}{lrrrrr}
\toprule
planner & median $\dv$ & induced & \textbf{induced} & scenarios & certified \\
        & (\si{\milli\meter\per\second}) & predicted & \textbf{measured} & affected & safe \\
\midrule
\texttt{no-maneuver}     &    0.0 &   0 & \textbf{0}   &  0/80 & --- \\
\texttt{legacy-lp}       & 1094.1 & 425 & \textbf{176} & 71/80 & \SI{0.0}{\percent} \\
\texttt{greedy-pairwise} &  899.2 & 229 & \textbf{97}  & 64/80 & \SI{0.0}{\percent} \\
\texttt{fuel-only}       &  715.7 & 200 & \textbf{71}  & 55/80 & \SI{100}{\percent} \\
\textbf{\texttt{fleet-safe}} & 867.7 & \textbf{0} & \textbf{0} & \textbf{0/80} & \SI{95.0}{\percent} \\
\bottomrule
\end{tabular}
\end{center}

The constraint drives its own predicted count from 200 to \textbf{zero}, and a
full SGP4 re-screen of the maneuvered catalog finds \textbf{zero} induced
conjunctions in every one of the eighty scenarios --- against 71 for the same
optimizer without the induced-conjunction rows, and 176 for the legacy
scalar-model planner. None of the uncoordinated planners passes certification,
and the one that spends the most fuel induces the most.

\paragraph{Across all nine families.} The table above isolates the family the
question is about. Repeating the sweep over all nine generator families --- 144
scenarios, \SI{2}{\meter\per\second} budget --- gives:

\begin{center}
\begin{tabular}{lrrrr}
\toprule
planner & median $\dv$ & induced & \textbf{induced} & scenarios \\
        & (\si{\milli\meter\per\second}) & predicted & \textbf{measured} & affected \\
\midrule
\texttt{no-maneuver}     &   0.0 &   0 & \textbf{0}  &   0/144 \\
\texttt{legacy-lp}       &  77.9 & 158 & \textbf{37} &  16/144 \\
\texttt{greedy-pairwise} &  12.5 &  73 & \textbf{18} &  13/144 \\
\texttt{fuel-only}       &  33.1 &  60 & \textbf{13} &  10/144 \\
\textbf{\texttt{lexicographic}} & 510.6 & 6 & \textbf{0} & \textbf{0/144} \\
\textbf{\texttt{fleet-safe}}    & 610.5 & 6 & \textbf{0} & \textbf{0/144} \\
\textbf{\texttt{milp-ops}}      & 513.1 & 6 & \textbf{0} & \textbf{0/144} \\
\bottomrule
\end{tabular}
\end{center}

Zero for every formulation carrying the induced-conjunction rows, across every
family --- not only the one the method was built around. The premium over all
families is defined on 43 of 144 scenarios, median \SI{0.00}{\percent}, p90
\SI{5.30}{\percent}, max \SI{15.97}{\percent}, \textbf{zero negative},
\SI{22.2}{\percent} with no zero-induced plan. The median is zero because most
families have no binding induced-conjunction row at all, which is the correct
answer for them; the \texttt{induced-cascade} distribution below is the one
carrying information.

Reaching zero required one structural correction, and it is worth stating
because the failure mode is not obvious. Placing rows at local minima of the
\emph{nominal} separation is not sufficient: for a co-orbital pair the nominal
separation is flat --- measured at \SI{9.345}{\kilo\meter} across every sampled
epoch --- so nominal minima are arbitrary, and the maneuver's own secular
along-track drift decides where the perturbed minimum falls. Plans therefore
satisfied every enumerated row, with a first-order model accurate to
\textbf{six metres}, while the true minimum sat two kilometres lower at an
epoch no row covered. Because evaluating the perturbed separation is cheap once
$x$ is known, the epochs that matter are discovered from the solution --- solve,
locate each pair's perturbed minimum, add one row, re-solve --- rather than
guessed beforehand. The cost is \SI{2.9}{\percent} more $\dv$ and no change to
the premium distribution.

The safety premium, measured at a fixed linearization with latent slack pinned
to zero and resolve slack capped at the fuel-only solution --- the construction
under which it is provably a lower bound:

\begin{center}
\begin{tabular}{lr}
\toprule
statistic & value \\
\midrule
defined & 70/80 \\
no zero-induced plan exists & 10/80 (\SI{12.5}{\percent}) \\
median & \textbf{\SI{4.41}{\percent}} \\
p90 & \textbf{\SI{16.00}{\percent}} \\
p99 & \SI{42.62}{\percent} \\
worst & \SI{54.69}{\percent} \\
zero-premium fraction & \SI{20}{\percent} \\
\textbf{negative (construction forbids it)} & \textbf{0} \\
\bottomrule
\end{tabular}
\end{center}

Stratified by how tightly the fleet is packed:

\begin{center}
\begin{tabular}{lrrrrr}
\toprule
neighbour spacing & $n$ & median & p90 & max & infeasible \\
\midrule
\SIrange{3}{5}{\kilo\meter}   &  9 & \textbf{\SI{29.67}{\percent}} & \SI{35.70}{\percent} & \SI{37.20}{\percent} & 6 \\
\SIrange{5}{7}{\kilo\meter}   & 12 & \SI{15.42}{\percent} & \SI{18.64}{\percent} & \SI{54.69}{\percent} & 0 \\
\SIrange{7}{9}{\kilo\meter}   & 15 & \SI{10.00}{\percent} & \SI{12.69}{\percent} & \SI{13.22}{\percent} & 0 \\
\SIrange{9}{12}{\kilo\meter}  & 21 & \SI{3.28}{\percent}  & \SI{6.10}{\percent}  & \SI{6.95}{\percent}  & 0 \\
\SIrange{12}{15}{\kilo\meter} & 23 & \textbf{\SI{0.00}{\percent}} & \SI{2.25}{\percent} & \SI{11.30}{\percent} & 4 \\
\bottomrule
\end{tabular}
\end{center}

Monotone across a factor of five in spacing and thirty in premium, with the
transition to infeasibility arriving below about \SI{5}{\kilo\meter}. The
median sits inside the \SIrange{1.1}{3.8}{\percent} range Pavanello et
al.~\cite{pavanello2024} report for a single spacecraft buying stronger
probability margins across five, seven and ten simultaneous encounters --- a
different formulation measuring a different thing, so corroboration rather than
confirmation, but the closest published quantity.

\paragraph{A sampled grid cannot locate a fast close approach, and the tidal
inflation assumes it can.} The result above is stated over a sweep that varies
the $\dv$ budget as well as the geometry, and that is deliberate: an earlier
version of it, swept at a single budget, was wrong. An independent verifier
given read-only access reproduced zero at \SI{2}{\meter\per\second} and then
broke it at \SI{0.5}{} and \SI{1}{\meter\per\second}, the latter being the
implementation's own default.

The failing pair crosses at \SI{7.24}{\kilo\meter\per\second} and so sweeps
\SI{217}{\kilo\meter} per \SI{30}{\second} step. Its true minimum,
\SI{0.1429}{\kilo\meter}, falls between two samples reading \SI{76.3}{} and
\SI{140.8}{\kilo\meter}: the sampled separation is monotone increasing straight
through a minimum three orders of magnitude below either neighbour. No
predicate on sampled separations --- $\arg\min$, a local-minimum scan, a
decrease-then-increase test --- can detect it. The first-order model was not
at fault; it agreed with SGP4 to a median of \SI{2.2}{\meter} over the window.

Nor does \eqref{eq:tidal} cover it, and the reason is worth isolating. The term
$\tfrac12 A h^2$ bounds the \emph{curvature} of the relative trajectory, which
is the governing quantity only when the relative velocity at the minimum is
near zero --- a co-orbital pair, which is precisely the geometry the
formulation was developed against. For a crossing, the governing quantity is
the linear term $\tfrac12 v_{\text{rel}} h$, the very term \eqref{eq:tidal} was
introduced to avoid, and at \SI{7.24}{\kilo\meter\per\second} it is
\SI{108}{\kilo\meter} against a \SI{2.88}{\kilo\meter} floor. Measured across
families, $\tfrac12 v_{\text{rel}} h$ exceeds the floor by a factor of
\num{20} to \num{53} everywhere --- so the grid was certifying nothing, in any
family. Zero had been \emph{measured}, not guaranteed.

\begin{center}
\begin{tabular}{lrrr}
\toprule
family & worst $\tfrac12 v_{\text{rel}} h$ & floor & ratio \\
\midrule
\texttt{isolated-pair}   & \SI{73.6}{\kilo\meter}  & \SI{3.68}{\kilo\meter} & $20\times$ \\
\texttt{chain}           & \SI{111.2}{\kilo\meter} & \SI{2.88}{\kilo\meter} & $39\times$ \\
\texttt{induced-cascade} & \SI{142.5}{\kilo\meter} & \SI{3.45}{\kilo\meter} & $41\times$ \\
\texttt{star}            & \SI{185.2}{\kilo\meter} & \SI{3.52}{\kilo\meter} & $\mathbf{53\times}$ \\
\bottomrule
\end{tabular}
\end{center}

The range rate repairs it. $\dot r = (\delta \cdot \dot\delta)/\lVert\delta\rVert$
changes sign at every minimum however deep, so brackets are found from a sign
change in $\dot r$ rather than by comparing separations, and each bracket is
refined by cubic Hermite interpolation of the relative position from its
endpoints and their derivatives. This requires the exact derivative of $\Phi$,
which is available in closed form, and makes the search independent of how far
the pair moves within one step. Root-finding on the range rate is of course
standard practice in conjunction \emph{screening}; what seems not to be
appreciated is that a maneuver-planning constraint set built on sampled
separations inherits none of that rigour, and that the tidal-inflation argument
routinely used to justify a coarse planning grid is conditional on a
relative-velocity assumption that is not usually stated, let alone checked.

A second verifier, given the strengthened claim and told to find a new axis of
attack, swept \num{560} configurations --- budgets from \SIrange{50}{10000}{\milli\meter\per\second},
seeds to \num{60}, burn-slot counts, refinement iterations, a one-axis control
model, the polyhedral cost model, two probability thresholds --- and broke it on
exactly one: a \SI{120}{\second} enumeration step. The search had not failed;
the pair was never offered to it. Candidate pairs were derived from the rows the
enumeration placed, and admission tested $\lVert d \rVert \le G$ \emph{at
samples}, so a pair whose whole approach falls between two samples has no sample
inside the gate, receives no row, and is never a candidate.

This locates where $\tfrac12 v_{\text{rel}} h$ belongs, and the placement is the
result. On the constraint \emph{floor} it is vacuous --- \SI{108}{\kilo\meter}
against a \SI{2.9}{\kilo\meter} exclusion radius is an empty feasible set, which
is what motivated \eqref{eq:tidal}. On the \emph{gate} it is nearly free:
widening admission admits a few more candidate pairs, costing enumeration time
and nothing else, since a pair that turns out not to bind makes no plan worse.
Admission therefore tests $\lVert d \rVert \le G + \tfrac12 v_{\text{rel}} h$
with $v_{\text{rel}}$ taken per pair rather than as a catalog-wide maximum, and
candidate membership is read from what the gate admitted rather than from what
received a row. The two are different questions, and conflating them is what
allowed a coarse grid to discard a \SI{7.2}{\kilo\meter\per\second} crossing
outright.

Across \num{900} plannable configurations --- nine families, sixteen seeds,
four budgets, three planners --- this takes the measured induced count from
\textbf{6} to \textbf{0}, and takes the number of those failures that reported
a \emph{safe} certificate from \textbf{6 of 6} to none. \num{593} of \num{899}
plans change; the median $\dv$ change is \SI{+0.000}{\percent} and the p90 is
\SI{+16.2}{\percent}, so the cost falls entirely on the cases that were wrong.

\paragraph{The enumeration step is a correctness parameter.} Sweeping it as a
third axis --- \num{3066} non-vacuous checks over five families, sixteen seeds,
four budgets, three steps and three planners, plus \num{366} more varying burn
slots, refinement iterations, control axes, cost model, probability threshold
and risk level --- locates the boundary:

\begin{center}
\begin{tabular}{lrr}
\toprule
enumeration step & measured induced & of which certified ``safe'' \\
\midrule
\textbf{\SI{30}{\second} (default)} & \textbf{0 / 900} & --- \\
\SI{60}{\second}  &  4 / 900 & 2 \\
\SI{120}{\second} & 24 / 900 & 17 \\
all other axes, at \SI{30}{\second} & \textbf{0 / 366} & --- \\
\bottomrule
\end{tabular}
\end{center}

Coarsening the step does not merely place fewer rows; it changes which
approaches are visible at all, since admission and placement both begin from
sampled separations. The other six parameters swept move the result not at all.
We report the guarantee as measured at \SI{30}{\second} and demonstrably
degrading above it, and the implementation flags a coarser grid as unvalidated
rather than letting it pass for the validated one. Three successive versions of
this result were each refuted on an axis the preceding sweep had not varied,
and all three root causes share one shape: a continuous quantity decided from
samples of it. For slow relative motion that is sound, which is why the
co-orbital geometry the formulation was developed against exposed none of them.

\paragraph{A pinned attainment level does not survive constraint generation.}
Lazy epoch generation interacts badly with one natural formulation, and the
interaction is worth stating because it is invisible until both are present.
A lexicographic safety-then-fuel planner solves twice: stage one prices slack
far above any achievable $\dv$ and reads off the total shortfall the geometry
allows; stage two pins that level and minimizes $\dv$ underneath it. With rows
generated lazily, stage one's level is measured against the \emph{enumerated}
rows. Stage two's refinement then discovers epochs those rows missed, adds a
row at each, and the added shortfall pushes the total above the pinned
level --- so stage two returns \texttt{infeasible} and the planner falls back
to the unrefined plan, which violates the very rows it just found. Across all
nine families this left exactly one SGP4-measured induced conjunction, in a
scenario where the linear model and SGP4 agreed to \SI{1.2}{\meter}: the model
was never wrong, the plan was simply never constrained at that epoch.
With stage one closed first, \texttt{lexicographic} reaches zero across all
144 scenarios, for \SI{41}{\percent} less $\dv$ than
\texttt{fleet-safe} --- it is entitled to stop at the best attainable safety
rather than insist on a feasible point.
\texttt{fleet-safe}, which pins nothing, absorbs the new rows and re-solves.
The fix is to close the epoch set on stage one before reading the level off it.
The general statement: \emph{an attainment level measured on one relaxation is
not transferable to a tighter one}, and with generated rows the final set is
known only after generation closes.

A related defect surfaced with it. The lazy routine grew the row set
internally but returned only the solution, so callers certified against the
originally enumerated rows --- and duly reported a plan
\emph{linearized-safe} while it violated a row the refinement itself had
added. The routine now returns the closed set. No plan changes; what the
certificate is entitled to claim does.

\paragraph{A hypothesis that did not survive, and its replacement.} We first
proposed a spectral predictor for the premium: the \emph{coupling number},
$\operatorname{cond}(BB^{\!\top})$ of the row-normalized projected
sensitivities, on the reasoning that ill-conditioning means several
conjunctions demand nearly collinear displacements. Its rank correlation with
the measured premium is \textbf{\num{0.074}} --- nothing, and on the
\texttt{induced-cascade} family it is constant.

What does predict it is the \emph{penetration index}: the violation the
\emph{fuel-optimal} plan commits on its worst induced-conjunction row,
normalized by the most any admissible plan could move that row,
\begin{equation}
\varrho \;=\; \max_p
  \frac{\big(f_p - \hat u_p^{\!\top}(d_p + B_p x_{\text{fuel}})\big)_+}
       {\sum_i D_i \,\lVert B_p \text{ restricted to satellite } i\rVert_\infty}.
\end{equation}
It is dimensionless, and reads as \emph{what fraction of the fleet's reach is
already spoken for by the cheapest plan's worst induced conjunction}. Measured
over 70 scenarios: Spearman \textbf{\num{0.989}}, univariate $R^2$ \num{0.866},
against \num{-0.855} for in-plane neighbour spacing and \num{0.462} for the
conflict dimension. A bake-off of thirteen further hand-constructed scalars,
and a random forest on sixty-six pre-screening features, found nothing better.

The sharper property is exact rather than statistical: $\varrho = 0$ \emph{if
and only if} the premium is zero, on \textbf{70 of 70} scenarios here and 71 of
71 in an independent run. If the fuel-optimal plan violates no
induced-conjunction row, coordination is free. We deliberately do not claim the
fitted slope as a calibrated magnitude --- it came out at \num{2.72} on one
dataset and \num{5.73} on another with a different budget, because it absorbs
the scenario's own $\dv$ scale. The index ranks, and it decides whether a
premium exists at all.


\subsection{The induced-risk frontier}

Traced at a fixed linearization with resolve shortfall pinned, over six
\texttt{induced-cascade} scenarios, Proposition~\ref{prop:frontier} holds to
machine precision: every curve non-increasing and convex, convexity residual
$\le \num{9e-19}$, slope--dual agreement $\le \num{1.5e-13}$, with breakpoints
recovered where the dual steps down. Premiums at $\varepsilon=0$ relative to
the unconstrained end ranged from \SI{0}{\percent} (spacing
\SI{10.6}{\kilo\meter}) to \SI{5.8}{\percent} (spacing \SI{6.9}{\kilo\meter}).

\subsection{Certified acceleration, and a negative result}
\label{sec:results-accel}

Before any of what follows, a profile is needed, because it changes which
number matters. On a 26-object scenario over four orbits the initial screen
costs \SI{14.7}{\second}, the SGP4 re-screen \SI{11.2}{\second}, latent-row
enumeration \SI{1.0}{\second}, and the linear program \SI{0.24}{\second}. The
LP is \textbf{1\%} of the runtime. Accelerating it, which is what the
predict-and-certify machinery does, is therefore worth at most one per cent
however well it works --- and the three changes that actually mattered were
exact algorithmic ones: memoising a rotation matrix per (object, epoch) and
writing out the cross products by hand (\textbf{2.1$\times$} on screening,
conjunction lists bit-identical), restricting the re-screen to pairs touching a
burn since nothing else moved (\textbf{2.5--14$\times$}, induced counts
identical), and the lazy closure below (up to \textbf{5$\times$} on the LP,
$\dv$ identical to 0.0e+00). End to end, \textbf{3.8$\times$}.

A fourth exact change removes the remaining scaling limit, which is memory
rather than time. The slack block is diagonal --- one slack column per resolve
or latent row --- so each latent row carries $n_{\text{lifted}}+1$ nonzeros
while the column count itself grows with the latent count, making dense storage
quadratic in the only dimension that grows. At \num{20840} rows the matrix is
\SI{0.23}{\percent} dense. Assembling coordinate triplets and materialising CSR
above four million elements takes peak assembly memory from \SI{3575}{\mega\byte}
to \SI{84}{\mega\byte}, a factor of \num{42.5}, with the matrix entrywise
identical, the solution vector bit-identical, and 56 of 56 benchmark rows
reproduced to \num{0.0e+00}.

With that said, the lazy-constraint closure returns the exact optimum for any
guess, and its speedup grows with the problem. With an oracle guess on the
\texttt{induced-cascade} family, varying the neighbour count and the latent
grid step:

\begin{center}
\begin{tabular}{rrrrr}
\toprule
latent rows & rows used & full solve & lazy solve & speedup \\
\midrule
   68 & 1 & \SI{4.5}{\milli\second}   & \SI{2.5}{\milli\second}  & $1.8\times$ \\
  358 & 2 & \SI{18.5}{\milli\second}  & \SI{5.4}{\milli\second}  & $3.5\times$ \\
  820 & 2 & \SI{48.5}{\milli\second}  & \SI{8.3}{\milli\second}  & $5.9\times$ \\
 1716 & 2 & \SI{145.8}{\milli\second} & \SI{16.1}{\milli\second} & $\mathbf{9.0\times}$ \\
\bottomrule
\end{tabular}
\end{center}

Objective gap exactly zero at every size. Exactness was verified against an
empty, a learned, and a deliberately adversarial guess over 27 problems: max
relative gap \num{1.4e-16}. An independent verifier repeated this with its own
seeds, spacings and neighbour counts over 20 further problems and 60
comparisons, and reported a maximum gap of \textbf{exactly zero}, bit-identical,
while confirming that the guess genuinely changes the number of rows used and
therefore the search path.

\paragraph{The learned predictor did not pay off.} A physics-informed graph
network (\num{892}k parameters) was trained on 320 scenarios across eight
families, split by scenario digest (258/39/23), with a loss combining
imitation of the solver's active set and duals with a physics term evaluating
real constraint residuals through the stored sensitivity tensors. Training is
exactly reproducible: two runs at the same seed gave identical validation loss
(difference \num{0}). Calibrated for recall, it retains \SI{64.6}{\percent} of
rows at \SI{100}{\percent} recall, or \SI{25.5}{\percent} at
\SI{96.2}{\percent}.

On 18 held-out large scenarios the objective gap remained exactly zero --- the
certification is indifferent to the model's quality --- but the median
end-to-end speedup including inference was $0.77\times$, against roughly
$4\times$ for simply starting from the \emph{empty} set. Two measurable
reasons: the active set is far smaller than the network's output (two binding
rows out of 1716, against 2--146 proposed), and inference costs
\SIrange{8}{56}{\milli\second} against a \SIrange{14}{180}{\milli\second}
solve.

A wider search confirmed and sharpened this. Four row-selection strategies were
compared, all returning a bit-identical objective: the empty guess
(\numrange{1.3}{10.6}$\times$, using 0--9 rows of 1716), the trained network
(\numrange{0.16}{1.91}$\times$, 119--199 rows), a \emph{certified} analytic
per-row prune (\numrange{0.65}{0.98}$\times$), and that prune combined with an
analytic a-posteriori guess (\numrange{0.43}{0.97}$\times$). The certified prune
deserves a note: a row can only be violated if the budget can move it past its
nominal margin, which is checkable in closed form, and over 8\,896 rows it lost
\textbf{zero} binding rows while pruning \SIrange{15}{65}{\percent}. It is
sound, and it is still slower than predicting nothing.

The claim that survives is about the mechanism: an active-set guess can be
turned into an exactly-certified speedup that grows with problem size. The claim
that does not survive is that \emph{any} predictor --- learned or analytic ---
supplies a better guess than the trivial one. The active set holds two to eleven
rows out of seventeen hundred, so a useful predictor would have to be cheaper
than free. We report this as measured, and note the general pattern it belongs
to: every acceleration that worked in this system replaced a guess with a cheap
way to \emph{know}.

\paragraph{Triage: a ceiling that means nothing.} A separate bake-off asked
whether a model could decide, from mean elements alone and before any
propagation, that a scenario needs no maneuver at all --- 53\% of a benchmark
sweep needs none and pays full screening cost anyway. Every model family
reached ROC-AUC \num{1.0} and specificity \num{1.0} at recall \num{1.0} on a
held-out split: a two-feature closed-form rule at \SI{0.57}{\micro\second},
logistic regression, gradient boosting, random forests, $k$-NN, MLPs, deep
ensembles, MC-dropout, Bayes-by-backprop, a Bayesian LSTM, DeepSets, a Set
Transformer and a 1D CNN.

That ceiling is an artifact. The label is constant within nine of the ten
scenario families, so the task is family recognition and two pooled statistics
accomplish it. A leave-one-family-out refit makes the failure explicit: with
\texttt{induced-cascade} held out, the rule skips \textbf{100 of 100} scenarios
that needed a maneuver --- recall \num{0.000}. Since the one intolerable error
for this system is a false negative, no triage gate is deployed. The genuinely
sound physics rule --- skip if no pair passes the perigee/apogee prefilter ---
fires on \textbf{0 of 423} scenarios, because every synthetic family lives on a
single \SI{550}{\kilo\meter} shell. A benchmark that could discriminate here
needs families whose label varies \emph{within} the family.



\section{Limitations}

Linearized relative motion, with second-order error scaling as
$\norm{\text{displacement}}^2 / (2\,m)$, reported with every claim. Low relative
velocity, where the TCA shift grows as $1/\norm{w}$ and the neglected relative
acceleration dominates: these events are flagged and reported separately.
Cross-track burns beyond about two orbits of lead, where differential $J_2$
reaches \SI{7.6}{\percent}. And every probability inherits TLE-grade
covariance, so no number here is operational.

The scenario families are synthetic and deterministic by design --- they exist
to produce known graph structure --- with one family replayed from real
catalog elements. A replication on operator-grade ephemerides would be the
natural next step, and would change the covariance story completely.

\section{Conclusion}

Induced conjunctions can be constrained inside a fleet maneuver optimizer
rather than chased by an outer loop, and the constraint set can be certified
complete under a $\dv$ budget. Doing so costs a median of a few percent extra
$\dv$ when the fleet is loosely packed, and becomes infeasible when it is
tightly packed --- a transition that is a property of the geometry, not of the
solver, and one that a soft-penalty formulation hides. The tradeoff that
remains, once the penalty is set above its exact threshold, is convex,
piecewise linear, and exactly computable.

\begin{thebibliography}{99}

\bibitem{armellin2021} R.~Armellin.
\newblock Collision avoidance maneuver optimization with a multiple-impulse convex formulation.
\newblock \emph{Acta Astronautica}, 186, 2021. arXiv:2101.07403.

\bibitem{pavanello2024} Z.~Pavanello, L.~Pirovano, R.~Armellin, A.~De~Vittori, P.~Di~Lizia.
\newblock Collision avoidance maneuvers optimization in the presence of multiple encounters.
\newblock arXiv:2406.03654, 2024.

\bibitem{chen2024} Y.~Chen et al.
\newblock Instability of self-driving satellite mega-constellation.
\newblock arXiv:2406.06068, 2024.

\bibitem{nasa2023} NASA Office of the Chief Engineer.
\newblock NASA spacecraft conjunction assessment and collision avoidance best practices handbook.
\newblock NASA/SP-20230002470 Rev.~1, 2023.

\bibitem{mueller2009} J.~B.~Mueller.
\newblock Onboard planning of collision avoidance maneuvers using robust optimization.
\newblock AIAA 2009-2051.

\bibitem{mao2017} Y.~Mao, M.~Szmuk, B.~A\c{c}{\i}kme\c{s}e.
\newblock Successive convexification of non-convex optimal control problems.
\newblock \emph{IFAC-PapersOnLine}, 50(1):4063--4069, 2017.

\bibitem{rivero2025} A.~S.~Rivero, R.~Vazquez, K.~Merz.
\newblock Advancing geometric conjunction filters to handle satellite manoeuvres.
\newblock 9th European Conference on Space Debris (SDC9), 2025.

\bibitem{hejduk2017} M.~D.~Hejduk, D.~E.~Pachura.
\newblock Conjunction assessment screening volume sizing.
\newblock NASA NTRS 20170007928.

\bibitem{klinkrad2005} H.~Klinkrad, J.~R.~Alarcon, N.~Sanchez.
\newblock Collision avoidance for operational ESA satellites.
\newblock ESA SP-587, 4th European Conference on Space Debris, 2005.

\bibitem{kuhl2025} W.~Kuhl, J.~Wang, D.~Eddy, M.~Kochenderfer.
\newblock Markov decision processes for satellite maneuver planning and collision avoidance.
\newblock arXiv:2501.02667, 2025.

\bibitem{bertsimas2019} D.~Bertsimas, B.~Stellato.
\newblock Online mixed-integer optimization in milliseconds.
\newblock arXiv:1907.02206, 2019.

\bibitem{cauligi2020} A.~Cauligi, P.~Culbertson, E.~Schmerling, M.~Pavone.
\newblock Learning mixed-integer convex optimization strategies for robot planning and control.
\newblock arXiv:2004.03736, 2020.

\bibitem{gasse2019} M.~Gasse, D.~Ch\'etalat, N.~Ferroni, L.~Charlin, A.~Lodi.
\newblock Exact combinatorial optimization with graph convolutional neural networks.
\newblock NeurIPS, 2019. arXiv:1906.01629.

\bibitem{uriot2020} T.~Uriot et al.
\newblock Spacecraft collision avoidance challenge: design and results of a machine learning competition.
\newblock arXiv:2008.03069, 2020.

\bibitem{bertsekas1999} D.~P.~Bertsekas.
\newblock \emph{Nonlinear Programming}. Athena Scientific, 2nd edition, 1999.

\bibitem{schrijver1986} A.~Schrijver.
\newblock \emph{Theory of Linear and Integer Programming}. Wiley, 1986.

\bibitem{alfano2005} S.~Alfano.
\newblock A numerical implementation of spherical object collision probability.
\newblock \emph{Journal of the Astronautical Sciences}, 53(1), 2005.

\end{thebibliography}

\end{document}
